\section{Introduction}
\vspace{1cm}
Compressed sensing is a signal processing technique that enables the reconstruction of a signal from a small number of measurements. Given the measurement $y$, the measurements matrix $A\in\mathbb{R}^{m\times n}$ and the noise $\eta\in\mathbb{R}^m$, the goal is to estimate the vector $x^*\in\mathbb{R}^n$ such that:
\begin{equation}
    y=Ax^*+\eta
\end{equation}
Since $m\ll n$, this system is underdetermined, then it is not possible to find a unique solution. The Nyquist--Shannon sampling theorem requires a certain amount of samples to perfectly reconstruct a signal, so compressed sensing exploits some features of the signal to address this problem in a context where the conditions of the aforementioned theorem are not met.
\\\\
The standard version of the algorithm assumes the signal to be $k$-sparse, meaning that it has only $k$ coefficients different from zero when represented in some domain. This assumption allows the signal retrieval through the use of algorithms for solving minimization problems, such as LASSO \cite{tibshirani1996} for the $L^1$ norm minimization, under conditions on $A$ established by the classical compressed sensing literature \cite{candes2006,donoho2006}.
\\\\
Following Bora et al. \cite{bora2017}, this experiment shifts from sparsity to structures derived from generative models, with the aim of assessing the goodness of these models when employed in compressed sensing. They feature a generator that maps vectors from  a low dimensional to a high dimensional space, so from $z\in\mathbb{R}^k$ to $G(z)\in\mathbb{R}^n$. By training the generator, with high probability it will output vectors similar to the ones belonging to the dataset it has been trained on.
\\\\
The idea behind the use of generative models in the context of compressed sensing is that we expect the vector $x^*$ to be close to some $G(z)$, which can be found with an optimization algorithm based on gradient descent, that will be later discussed in detail.